% questions.tex —— 2008 年·普通高等学校招生考试·新课标卷（琼、宁卷）·文科 · 题目内容片段（无答案）
% 由 tools/generate_projects.py 依来源总账自 2008/2008_数学真题_新课标文_试题.md 生成，勿手改（改源 MD 后重跑）


\gkpart{一、选择题}

\begin{question}
\gkti{1}已知集合 $M=\{x\mid(x+2)(x-1)<0\}$，$N=\{x\mid x+1<0\}$，则 $M\cap N=$（　）

（A）$(-1,1)$

（B）$(-2,1)$

（C）$(-2,-1)$

（D）$(1,2)$

\end{question}
\begin{question}
\gkti{2}双曲线 $\dfrac{x^2}{10}-\dfrac{y^2}{2}=1$ 的焦距为（　）

（A）$3\sqrt2$

（B）$4\sqrt2$

（C）$3\sqrt3$

（D）$4\sqrt3$

\end{question}
\begin{question}
\gkti{3}已知复数 $z=1-\mathrm i$，则 $\dfrac{z^2}{z-1}=$（　）

（A）$2$

（B）$-2$

（C）$2\mathrm i$

（D）$-2\mathrm i$

\end{question}
\begin{question}
\gkti{4}设 $f(x)=x\ln x$，若 $f'(x_0)=2$，则 $x_0=$（　）

（A）$\mathrm e^2$

（B）$\mathrm e$

（C）$\dfrac{\ln2}2$

（D）$\ln2$

\end{question}
\begin{question}
\gkti{5}已知平面向量 $\boldsymbol a=(1,-3)$，$\boldsymbol b=(4,-2)$，$\lambda\boldsymbol a+\boldsymbol b$ 与 $\boldsymbol a$ 垂直，则 $\lambda=$（　）

（A）$-1$

（B）$1$

（C）$-2$

（D）$2$

\Needspace{0.4\textheight}
\end{question}
\begin{question}
\gkti{6}下面的程序框图，如果输入三个实数 $a$，$b$，$c$，要求输出这三个数中最大的数，那么在空白的判断框中，应该填入下面四个选项中的（　）

（A）$c>x$

（B）$x>c$

（C）$c>b$

（D）$b>c$
\nopagebreak[4]
\begin{center}% 2008 年新课标卷文科第 6 题：输出三个数中最大值的程序框图（第二个判断框留空）。
% 来源 PDF：gaokaomath-src/普通高考/2008/2008新课标文(海南,宁夏).pdf
% SHA-256：853bce8a468fcc18b76717db6aa6fc777781a4da61ecde6120c1c4cae51cebff
% 原页：2008新课标文(海南,宁夏).pdf 第 1 页（A3 三栏）第 1 栏，书页码 469。
% 绘制依据：按原页流程：开始 → 输入 a,b,c → x←a → 判断 b>x（是则 x←b）→ 判断框（原页留空）
%   （是则 x←c）→ 输出 x → 结束。
\begin{tikzpicture}[node distance=0.42cm, semithick, every node/.style={inner sep=2.5pt, font=\small},
  box/.style={draw, minimum width=1.9cm, minimum height=0.5cm},
  io/.style={draw, trapezium, trapezium left angle=72, trapezium right angle=108,
             minimum width=2.2cm, minimum height=0.5cm},
  dec/.style={draw, diamond, aspect=2.2, minimum width=1.9cm, minimum height=0.75cm}]

  \node[box, rounded corners=2pt] (start) {开始};
  \node[io, below=of start] (input) {输入 $a,b,c$};
  \node[box, below=of input] (xa) {$x\leftarrow a$};
  \node[dec, below=0.45cm of xa] (d1) {$b>x$};
  \node[box, right=1.5cm of d1] (xb) {$x\leftarrow b$};
  \node[dec, below=0.75cm of d1, minimum width=1.9cm, minimum height=0.75cm] (d2) {};
  \node[box, right=1.5cm of d2] (xc) {$x\leftarrow c$};
  \node[io, below=0.7cm of d2] (out) {输出 $x$};
  \node[box, rounded corners=2pt, below=of out] (end) {结束};

  \draw[->,>=stealth] (start) -- (input);
  \draw[->,>=stealth] (input) -- (xa);
  \draw[->,>=stealth] (xa) -- (d1);
  \draw[->,>=stealth] (d1) -- node[above, font=\scriptsize] {是} (xb);
  \draw (xb) |- ($(d1.south)!0.5!(d2.north)$);
  \draw[->,>=stealth] (d1) -- node[left=4pt, font=\scriptsize] {否} (d2);
  \draw[->,>=stealth] (d2) -- node[above, font=\scriptsize] {是} (xc);
  \draw (xc) |- ($(d2.south)!0.5!(out.north)$);
  \draw[->,>=stealth] (d2) -- node[left=4pt, font=\scriptsize] {否} (out);
  \draw[->,>=stealth] (out) -- (end);
\end{tikzpicture}
\end{center}
\nopagebreak[4]
\end{question}
\begin{question}
\gkti{7}已知 $a_1>a_2>a_3>0$，则使得 $(1-a_ix)^2<1$（$i=1,2,3$）都成立的 $x$ 取值范围是（　）

（A）$\left(0,\dfrac1{a_1}\right)$

（B）$\left(0,\dfrac2{a_1}\right)$

（C）$\left(0,\dfrac1{a_3}\right)$

（D）$\left(0,\dfrac2{a_3}\right)$

\end{question}
\begin{question}
\gkti{8}设等比数列 $\{a_n\}$ 的公比 $q=2$，前 $n$ 项和为 $S_n$，则 $\dfrac{S_4}{a_2}=$（　）

（A）$2$

（B）$4$

（C）$\dfrac{15}2$

（D）$\dfrac{17}2$

\end{question}
\begin{question}
\gkti{9}平面向量 $\boldsymbol a$，$\boldsymbol b$ 共线的充要条件是（　）

（A）$\boldsymbol a$，$\boldsymbol b$ 方向相同

（B）$\boldsymbol a$，$\boldsymbol b$ 两向量中至少有一个为零向量

（C）$\exists\lambda\in\mathbf R$，$\boldsymbol b=\lambda\boldsymbol a$

（D）存在不全为零的实数 $\lambda_1$，$\lambda_2$，$\lambda_1\boldsymbol a+\lambda_2\boldsymbol b=\boldsymbol 0$

\end{question}
\begin{question}
\gkti{10}点 $P(x,y)$ 在直线 $4x+3y=0$ 上，且 $x$，$y$ 满足 $-14\leqslant x-y\leqslant7$，则点 $P$ 到坐标原点距离的取值范围是（　）

（A）$[0,5]$

（B）$[0,10]$

（C）$[5,10]$

（D）$[5,15]$

\end{question}
\begin{question}
\gkti{11}函数 $f(x)=\cos2x+2\sin x$ 的最小值和最大值分别为（　）

（A）$-1$，$1$

（B）$-2$，$2$

（C）$-3$，$\dfrac32$

（D）$-2$，$\dfrac32$

\end{question}
\begin{question}
\gkti{12}已知平面 $\alpha\perp$ 平面 $\beta$，$\alpha\cap\beta=l$，点 $A\in\alpha$，$A\notin l$，直线 $AB\parallel l$，直线 $AC\perp l$，直线 $m\parallel\alpha$，$m\parallel\beta$，则下列四种位置关系中，不一定成立的是（　）

（A）$AB\parallel m$

（B）$AC\perp m$

（C）$AB\parallel\beta$

（D）$AC\perp\beta$

\end{question}
\gkpart{二、填空题}

\begin{question}
\gkti{13}已知 $\{a_n\}$ 为等差数列，$a_3+a_8=22$，$a_6=7$，则 $a_5=$\gkblank{}。

\end{question}
\begin{question}
\gkti{14}一个六棱柱的底面是正六边形，其侧棱垂直底面。已知该六棱柱的顶点都在同一个球面上，且该六棱柱的高为 $\sqrt3$，底面周长为 $3$，则这个球的体积为\gkblank{}。

\end{question}
\begin{question}
\gkti{15}过椭圆 $\dfrac{x^2}5+\dfrac{y^2}4=1$ 的右焦点作一条斜率为 $2$ 的直线与椭圆交于 $A$、$B$ 两点，$O$ 为坐标原点，则 $\triangle OAB$ 的面积为\gkblank{}。

\end{question}
\begin{question}
\gkti{16}从甲、乙两品种的棉花中各抽测了 $25$ 根棉花的纤维长度（单位：mm），结果如下：

甲品种：$271$，$273$，$280$，$285$，$285$，$287$，$292$，$294$，$295$，$301$，$303$，$303$，$307$，$308$，$310$，$314$，$319$，$323$，$325$，$325$，$328$，$331$，$334$，$337$，$352$；

乙品种：$284$，$292$，$295$，$304$，$306$，$307$，$312$，$313$，$315$，$315$，$316$，$318$，$318$，$320$，$322$，$322$，$324$，$327$，$329$，$331$，$333$，$336$，$337$，$343$，$356$。

由以上数据设计了如下茎叶图：

\begin{center}
\begin{tabular}{rcl}
\toprule
甲 & 茎 & 乙 \\
\midrule
$3\ 1$ & $27$ &  \\
$7\ 5\ 5\ 0$ & $28$ & $4$ \\
$5\ 4\ 2$ & $29$ & $2\ 5$ \\
$8\ 7\ 3\ 3\ 1$ & $30$ & $4\ 6\ 7$ \\
$9\ 4\ 0$ & $31$ & $2\ 3\ 5\ 5\ 6\ 8\ 8$ \\
$8\ 5\ 5\ 3$ & $32$ & $0\ 2\ 2\ 4\ 7\ 9$ \\
$7\ 4\ 1$ & $33$ & $1\ 3\ 6\ 7$ \\
 & $34$ & $3$ \\
$2$ & $35$ & $6$ \\
\bottomrule
\end{tabular}
\end{center}

根据以上茎叶图，对甲、乙两品种棉花的纤维长度作比较，写出两个统计结论：

① \gkblank{}；

② \gkblank{}。

\end{question}
\gkpart{三、解答题}

\Needspace{0.4\textheight}
\begin{question}
\gkti{17}如图，$\triangle ACD$ 是等边三角形，$\triangle ABC$ 是等腰直角三角形，$\angle ACB=90^\circ$，$BD$ 交 $AC$ 于 $E$，$AB=2$。

（1）求 $\cos\angle CBE$ 的值；

（2）求 $AE$。
\nopagebreak[4]
\begin{center}% 2008 年新课标卷文科第 17 题：△ACD 为等边三角形，△ABC 为等腰直角三角形（∠ACB=90°），
%   BD 交 AC 于 E，AB=2。
% 来源 PDF：gaokaomath-src/普通高考/2008/2008新课标文(海南,宁夏).pdf
% SHA-256：853bce8a468fcc18b76717db6aa6fc777781a4da61ecde6120c1c4cae51cebff
% 原页：2008新课标文(海南,宁夏).pdf 第 1 页（A3 三栏）第 3 栏，书页码 469。
% 绘制依据：取 A(0,0)、B(2,0)、C(1,1)（∠ACB=90°、CA=CB=√2），
%   等边三角形 ACD 的第三个顶点 D 位于 AC 的另一侧，D=(-0.366,1.366)；E 为 BD 与 AC 的交点。
\begin{tikzpicture}[scale=2.0, semithick, every node/.style={inner sep=1.4pt}]
  \coordinate (A) at (0.00,0.00);
  \coordinate (B) at (2.00,0.00);
  \coordinate (C) at (1.00,1.00);
  \coordinate (D) at (-0.366,1.366);
  \coordinate (E) at (0.732,0.732);

  \draw (A) -- (B) -- (C) -- cycle;
  \draw (A) -- (D) -- (C);
  \draw (B) -- (D);

  % 直角标记
  \draw ($(C)!0.12!(A)$) -- ($(C)!0.12!(A)+($(C)!0.12!(B)-(C)$)$) -- ($(C)!0.12!(B)$);

  \fill (E) circle (0.7pt);

  \node[below left] at (A) {$A$};
  \node[below right] at (B) {$B$};
  \node[above right=0.5pt] at (C) {$C$};
  \node[left] at (D) {$D$};
  \node[below right=0.5pt] at (E) {$E$};
\end{tikzpicture}
\end{center}
\nopagebreak[4]
\Needspace{0.4\textheight}
\end{question}
\begin{question}
\gkti{18}如下的三个图中，上面的是一个长方体截去一个角所得多面体的直观图，它的正视图和侧视图在下面画出（单位：cm）。

（1）在正视图下面，按照画三视图的要求画出该多面体的俯视图；

（2）按照给出的尺寸，求该多面体的体积；

（3）在所给直观图中连结 $BC'$，证明：$BC'\parallel$ 面 $EFG$。
\nopagebreak[4]
\begin{center}% 2008 年新课标卷文科第 18 题：长方体截去一个角所得多面体的直观图，及它的正视图、侧视图
%   （单位：cm）。长方体 6×4×4，顶点 A' 处沿三条棱各截去 2。
% 来源 PDF：gaokaomath-src/普通高考/2008/2008新课标文(海南,宁夏).pdf
% SHA-256：853bce8a468fcc18b76717db6aa6fc777781a4da61ecde6120c1c4cae51cebff
% 原页：2008新课标文(海南,宁夏).pdf 第 1 页（A3 三栏）第 3 栏，书页码 469。
% 绘制依据：A(0,0,0)、B(6,0,0)、C(6,4,0)、D(0,4,0)、A'(0,0,4)、B'(6,0,4)、C'(6,4,4)、D'(0,4,4)，
%   截口点 E(2,0,4)、F(0,2,4)、G(0,0,2)；斜投影 (x,y,z)->(x+0.6y, z+0.35y)，取比例 0.42。
%   可见棱按「前面五边形 + 上底五边形 + 右侧面」联结，被遮挡的 A—D、D—C、D—D' 用虚线。
\begin{tikzpicture}[scale=1.0, semithick, every node/.style={inner sep=1.2pt, font=\small}]

  % ---------- 直观图 ----------
  \begin{scope}[shift={(0,0)}, scale=1.6]
    \coordinate (A)  at (0.00,0.00);
    \coordinate (B)  at (2.52,0.00);
    \coordinate (C)  at (3.528,0.588);
    \coordinate (D)  at (1.008,0.588);
    \coordinate (Bp) at (2.52,1.68);
    \coordinate (Cp) at (3.528,2.268);
    \coordinate (Dp) at (1.008,2.268);
    \coordinate (E)  at (0.84,1.68);
    \coordinate (F)  at (0.504,1.974);
    \coordinate (G)  at (0.00,0.84);

    % 前面五边形 A-B-B'-E-G
    \draw (A) -- (B) -- (Bp) -- (E) -- (G) -- cycle;
    % 右侧面
    \draw (B) -- (C) -- (Cp) -- (Bp);
    % 上底五边形 B'-C'-D'-F-E
    \draw (Bp) -- (Cp) -- (Dp) -- (F) -- (E);
    % 截口三角形
    \draw (E) -- (F) -- (G);

    % 被遮挡的棱
    \draw[dashed] (A) -- (D);
    \draw[dashed] (D) -- (C);
    \draw[dashed] (D) -- (Dp);

    \node[below left] at (A) {$A$};
    \node[below] at (B) {$B$};
    \node[right] at (C) {$C$};
    \node[left] at (D) {$D$};
    \node[above left] at (Dp) {$D'$};
    \node[right] at (Cp) {$C'$};
    \node[right] at (Bp) {$B'$};
    \node[below right=0.5pt] at (E) {$E$};
    \node[right=1pt] at (F) {$F$};
    \node[left=1pt] at (G) {$G$};
  \end{scope}

  % ---------- 正视图 ----------
  \begin{scope}[shift={(0.5,-3.9)}, scale=0.34]
    \draw (0,0) rectangle (6,4);
    \draw (0,2) -- (2,4);
    \draw[<->,>=stealth] (0,5.2) -- (6,5.2) node[midway, above, font=\scriptsize] {$6$};
    \draw[<->,>=stealth] (0,4.4) -- (2,4.4) node[midway, above, font=\scriptsize] {$2$};
    \draw[<->,>=stealth] (6.7,0) -- (6.7,4) node[midway, right, font=\scriptsize] {$4$};
    \node[below, font=\small] at (3,-0.8) {（正视图）};
  \end{scope}

  % ---------- 侧视图 ----------
  \begin{scope}[shift={(4.3,-3.9)}, scale=0.34]
    \draw (0,0) rectangle (4,4);
    \draw (2,4) -- (4,2);
    \draw[<->,>=stealth] (0,-0.8) -- (4,-0.8) node[midway, below, font=\scriptsize] {$4$};
    \draw[<->,>=stealth] (0,4.8) -- (2,4.8) node[midway, above, font=\scriptsize] {$2$};
    \draw[<->,>=stealth] (4.7,4) -- (4.7,2) node[midway, right, font=\scriptsize] {$2$};
    \node[below, font=\small] at (2,-1.5) {（侧视图）};
  \end{scope}
\end{tikzpicture}
\end{center}
\nopagebreak[4]
\end{question}
\begin{question}
\gkti{19}为了了解《中华人民共和国道路交通安全法》在学生中的普及情况，调查部门对某校 $6$ 名学生进行问卷调查。$6$ 人得分情况如下：$5$，$6$，$7$，$8$，$9$，$10$。把这 $6$ 名学生的得分看成一个总体。

（1）求该总体的平均数；

（2）用简单随机抽样方法从这 $6$ 名学生中抽取 $2$ 名，他们的得分组成一个样本。求该样本平均数与总体平均数之差的绝对值不超过 $0.5$ 的概率。

\end{question}
\begin{question}
\gkti{20}已知 $m\in\mathbf R$，直线 $l:mx-(m^2+1)y=4m$ 和圆 $C:x^2+y^2-8x+4y+16=0$。

（1）求直线 $l$ 斜率的取值范围；

（2）直线 $l$ 能否将圆 $C$ 分割成弧长的比值为 $\dfrac12$ 的两段圆弧？为什么？

\end{question}
\begin{question}
\gkti{21}设函数 $f(x)=ax-\dfrac bx$，曲线 $y=f(x)$ 在点 $(2,f(2))$ 处的切线方程为 $7x-4y-12=0$。

（1）求 $f(x)$ 的解析式；

（2）证明：曲线 $y=f(x)$ 上任一点处的切线与直线 $x=0$ 和直线 $y=x$ 所围成的三角形面积为定值，并求此定值。

\end{question}
\gkpart{四、选做题}

\Needspace{0.4\textheight}
\begin{question}
\gkti{22}如图，过圆 $O$ 外一点 $M$ 作它的一条切线，切点为 $A$，过 $A$ 点作直线 $AP$ 垂直直线 $OM$，垂足为 $P$。

（1）证明：$OM\cdot OP=OA^2$；

（2）$N$ 为线段 $AP$ 上一点，直线 $NB$ 垂直直线 $ON$，且交圆 $O$ 于 $B$ 点。过 $B$ 点的切线交直线 $ON$ 于 $K$。证明：$\angle OKM=90^\circ$。
\nopagebreak[4]
\begin{center}% 2008 年新课标卷文科第 22 题：过圆 O 外一点 M 作切线 MA，过 A 作 AP ⊥ OM 于 P；
%   N 在 AP 上，NB ⊥ ON 交圆于 B，过 B 的切线交 ON 于 K。
% 来源 PDF：gaokaomath-src/普通高考/2008/2008新课标文(海南,宁夏).pdf
% SHA-256：853bce8a468fcc18b76717db6aa6fc777781a4da61ecde6120c1c4cae51cebff
% 原页：2008新课标文(海南,宁夏).pdf 第 2 页（A3 三栏）第 3 栏，书页码 470。
% 绘制依据：取单位圆 O(0,0)、M(2,0)、A(1/2,√3/2)、P(1/2,0)、N(1/2,1/2)、B(0,1)、K(1,1)。
\begin{tikzpicture}[scale=2.4, semithick, every node/.style={inner sep=1.3pt}]
  \coordinate (O) at (0,0);
  \coordinate (M) at (2,0);
  \coordinate (A) at (0.5,0.866);
  \coordinate (P) at (0.5,0);
  \coordinate (N) at (0.5,0.5);
  \coordinate (B) at (0,1);
  \coordinate (K) at (1,1);

  \draw (O) circle [radius=1];

  % 半径 OA、切线 MA 与 OM
  \draw (O) -- (A);
  \draw (O) -- (M);
  \draw (M) -- (A);

  % AP、ON 及其延长到 K
  \draw (A) -- (P);
  \draw (O) -- (K);

  % NB 与过 B 的切线 BK
  \draw (N) -- (B);
  \draw (B) -- (K);

  % 直角标记：∠OAM 与 ∠OKM
  \draw ($(A)!0.1!(O)$) -- ($(A)!0.1!(O)+($(A)!0.1!(M)-(A)$)$) -- ($(A)!0.1!(M)$);
  \draw ($(K)!0.12!(O)$) -- ($(K)!0.12!(O)+($(K)!0.12!(M)-(K)$)$) -- ($(K)!0.12!(M)$);

  \fill (P) circle (0.9pt);
  \fill (N) circle (0.9pt);

  \node[below left] at (O) {$O$};
  \node[below] at (M) {$M$};
  \node[anchor=west] at ($(A)+(0.14,0.03)$) {$A$};
  \node[below] at (P) {$P$};
  \node[anchor=west] at ($(N)+(0.13,-0.02)$) {$N$};
  \node[anchor=south] at ($(B)+(0,0.12)$) {$B$};
  \node[anchor=south west] at ($(K)+(0.04,0.06)$) {$K$};
\end{tikzpicture}
\end{center}
\nopagebreak[4]
\end{question}
\begin{question}
\gkti{23}已知曲线 $C_1:\begin{cases}x=\cos\theta,\\ y=\sin\theta,\end{cases}$（$\theta$ 为参数），曲线 $C_2:\begin{cases}x=\dfrac{\sqrt2}2t-\sqrt2,\\ y=\dfrac{\sqrt2}2t,\end{cases}$（$t$ 为参数）。

（1）指出 $C_1$，$C_2$ 各是什么曲线，并说明 $C_1$ 与 $C_2$ 公共点的个数；

（2）若把 $C_1$，$C_2$ 上各点的纵坐标都压缩为原来的一半，分别得到曲线 $C_1'$，$C_2'$。写出 $C_1'$，$C_2'$ 的参数方程。$C_1'$ 与 $C_2'$ 公共点的个数和 $C_1$ 与 $C_2$ 公共点的个数是否相同？说明你的理由。

\end{question}